\hypertarget{proof}{%
\section{Green's Problem 64}\label{proof}}

\textbf{Theorem.} There are infinitely many primes \(p\) such that
\(\Omega(p-2)\) is odd, where \(\Omega(n)\) counts prime factors with
multiplicity.

This answers Problem 64 in
\href{https://people.maths.ox.ac.uk/greenbj/papers/open-problems.pdf}{Ben
Green's list}, p.~31. We prove the assertion for the fixed shift
\(p-2\). Throughout the argument, \(\lambda(n)=(-1)^{\Omega(n)}\) for
positive integers \(n\).

The proof constructs a nonnegative weight supported on integers \(u\)
with \(\lambda(u)=-1\), and shows that its sum over prime values of
\(u+2\) is positive in every sufficiently large multiplicative interval.
The central analytic input is the shifted-correlation estimate (B18),
which gives the Type II estimate (B3) uniformly in its bounded companion
coefficient. Its proof transfers cancellation on independent prime
coordinates to arithmetic histories, retaining the signed comparison and
accounting for repeated labels. A weighted sieve then removes most
composites; a separate estimate bounds those whose two prime factors are
close in size. The constants in the bound for the remaining composites
are chosen before the width of the balanced range; this permits that
last contribution to be subtracted.

The argument is divided into \protect\hyperlink{weights}{weights and
their mass}, \protect\hyperlink{correlation}{correlation and bilinear
estimates}, and the \protect\hyperlink{sieve}{sieve estimates}. The
present file defines the objects used in the conclusion and derives the
theorem from the two quantitative sieve estimates.

\Needspace{6\baselineskip}

\hypertarget{proof-weights-with-a-prescribed-parity}{%
\subsection{Weights with a prescribed
parity}\label{proof-weights-with-a-prescribed-parity}}

Write \(x=e^L\), \(T=\log L\), and suppose \(L\) is sufficiently large.
Set

\[
r=20,\qquad \sigma=\frac1{44},\qquad s_j=\sigma 2^{-j}L,
\qquad d=\max\{j\ge0:s_j\ge20T\}.
\]

For \(0\le j\le d\) let

\[
Q_j=\{q\text{ prime}:s_j\le\log q\le\tfrac65s_j\}.
\]

Also take every prime interval

\[
P_j=\{p\text{ prime}:\tfrac32s_j\le\log p\le\tfrac95s_j\}
\]

whose logarithmic interval is contained in either \([L^{1/10},L^{1/5}]\)
or \([L^{3/10},L^{2/5}]\). Let \(\mathcal G\) be this set of indices and
\(\mathcal P=\bigcup_{g\in\mathcal G}P_g\). All the \(Q_j\) and \(P_g\)
are pairwise disjoint. Put

\[
V_j=\sum_{q\in Q_j}\frac1q,\qquad M_g=\sum_{p\in P_g}\frac1p.
\]

The bounds used below hold uniformly for large \(L\):

\[
\frac1{12}\le V_j\le4,\qquad
\nu\le M_g\le4,\qquad \nu=\frac14\log\frac{12}{11}>0.
\]

For a positive integer \(u\), separate its prime factors as
\(u=u_*u_{\mathcal P}\), where \(u_*\) has no prime factor in
\(\mathcal P\) and \(u_{\mathcal P}\) has no prime factor outside
\(\mathcal P\). Define

\[
\omega_g(u)=\#\{p\in P_g:p\mid u\},\qquad
W(u)=\prod_{g\in\mathcal G}
       \frac{\omega_g(u)2^{-(\omega_g(u)-1)}}{M_g},
\]

with a zero factor when \(\omega_g(u)=0\). Notice the distinction
between \(\omega_g\), which counts distinct primes, and \(\Omega\),
which counts multiplicities.

Let \(a(v)\) be zero unless all prime factors of \(v\) lie in the
\(Q_j\) and \(v\) has exactly \(20\) factors, with multiplicity, in each
\(Q_j\). On that support set

\[
a(v)=\prod_{q\mid v}\frac1{v_q(v)!}.
\]

Equivalently, \(a(v)\) counts ordered choices of twenty primes in each
\(Q_j\) with product \(v\), divided by \((20!)^{d+1}\). Repeated primes
are allowed. Define

\[
E(u)=(-1)^{\sum_{p\in\mathcal P}v_p(u)},\qquad
A_\rho(u)=W(u)a(u_*)\frac{1+\rho E(u)}2
\quad(\rho\in\{-1,1\}).
\]

We write \(A_0(u)=W(u)a(u_*)\) for the unsigned weight; the subscript
zero is a separate notation, outside the two-sign definition above.

Since \(20(d+1)\) is even, \(a(u_*)>0\) implies \(\lambda(u_*)=1\).
Consequently

\[
A_\rho(u)>0\quad\Longrightarrow\quad\lambda(u)=\rho. \tag{1}
\]

This identity uses prime multiplicities throughout; it does not require
\(u\) to be squarefree.

Fix a smooth function \(\Psi:\mathbb R\to[0,1]\), supported in \([1,2]\)
and positive at some point of \((1,2)\). Such a function is obtained by
translating and rescaling a smooth bump function. Set

\[
w_\rho(u)=A_\rho(u)\Psi(u/x),\qquad
X_\rho=\sum_{u\ge1}w_\rho(u).
\]

Set \(w_\rho(0)=0\), consistently with \(\Psi(0)=0\). All these sums are
finite because \(w_\rho(u)\) vanishes outside \([x,2x]\). The mass
estimate proved in
\protect\hyperlink{weights-localization}{Section~\ref{weights-localization}}
gives a constant \(c_0>0\), depending on \(\Psi\) but not on \(L\) or
\(\rho\), such that

\[
X_\rho\ge c_0xL^{-221}>0 \tag{2}
\]

for both signs and all sufficiently large \(L\).

\Needspace{6\baselineskip}

\hypertarget{proof-the-two-sieve-estimates}{%
\subsection{The two sieve
estimates}\label{proof-the-two-sieve-estimates}}

Let

\[
\mathfrak S=2\prod_{p\ge3,\ p\text{ prime}}
                  \left(1-\frac1{(p-1)^2}\right).
\]

The product converges, and \(\mathfrak S\ge1\); the elementary product
bound is proved in
\protect\hyperlink{sieve-local-densities}{Section~\ref{sieve-local-densities}}.
Write

\[
\mathcal M_\rho(L)=\sum_{u+2\text{ prime}}w_\rho(u)
\]

and, for \(0<\kappa<1/50\), write

\[
\mathcal C_\rho(L,\kappa)
 =\sum_{\substack{u+2\text{ composite}\\
                   P^-(u+2)>x^{1/2-\kappa}}}w_\rho(u),
\]

where \(P^-(n)\) is the least prime factor of \(n>1\).

The weighted sieve and the subtraction of composites outside the
balanced range give, for every fixed \(0<\kappa<1/50\) and all
sufficiently large \(L\),

\[
\mathcal M_{-1}(L)+\mathcal C_{-1}(L,\kappa)
 \ge\frac{\mathfrak S}{2}\frac{X_{-1}}L. \tag{3}
\]

The negative sign is the one needed here. The complete subtraction,
including the order of the approximation parameters, is given in
\protect\hyperlink{sieve-removing-the-unbalanced-composites}{Section~\ref{sieve-removing-the-unbalanced-composites}}.

The bound for the remaining composites has the following uniform form:
there exists \(C_{\mathrm{bal}}>0\) such that, for every
\(0<\kappa<1/50\) and every \(\varepsilon>0\), both signs satisfy

\[
\mathcal C_\rho(L,\kappa)
 \le (C_{\mathrm{bal}}\kappa+\varepsilon)\frac{X_\rho}L
 \tag{4}
\]

once \(L\) is sufficiently large. The threshold may depend on \(\kappa\)
and \(\varepsilon\); \(C_{\mathrm{bal}}\) does not. Its proof in
\protect\hyperlink{sieve-balanced-composites}{Section~\ref{sieve-balanced-composites}}
uses an upper bound for positive factor pairs in a rectangle, followed
by a count of compatible dyadic rectangles. It includes pairs with equal
prime factors.

\Needspace{6\baselineskip}

\hypertarget{proof-positive-prime-mass-and-infinitude}{%
\subsection{Positive prime mass and
infinitude}\label{proof-positive-prime-mass-and-infinitude}}

Choose

\[
\kappa=\min\left\{\frac1{100},
                  \frac{\mathfrak S}{8C_{\mathrm{bal}}}\right\},
\qquad \varepsilon=\frac{\mathfrak S}{8}.
\]

Then \(0<\kappa<1/50\) and
\(C_{\mathrm{bal}}\kappa+\varepsilon\le\mathfrak S/4\). For this fixed
choice, intersect the eventual ranges in (2), (3), and (4). Subtracting
(4) from (3) gives

\[
\mathcal M_{-1}(L)
 \ge\frac{\mathfrak S}{4}\frac{X_{-1}}L>0. \tag{5}
\]

Every summand in \(\mathcal M_{-1}\) is nonnegative. Therefore some
positive integer \(u\) satisfies

\[
u+2\text{ is prime},\qquad A_{-1}(u)>0,\qquad\Psi(u/x)>0.
\]

The support of \(\Psi\) gives \(x\le u\le2x\). By (1),
\(\lambda(u)=-1\), so \(\Omega(u)\) is odd. For \(p=u+2\) we have
\(p-2=u>0\) and hence \(\Omega(p-2)\) odd.

Given any natural number \(M\), choose \(L\) in the eventual range of
(5) with \(e^L>M\). The resulting prime satisfies \(p=u+2>M\). Thus the
set of primes in the theorem is unbounded, and therefore infinite.
\(\square\)
